Este trabalho investigou se a substituição do \textit{embedding} monolítico do DGI por representações desacopladas e especializadas poderia produzir uma base mais adequada para as tarefas de classificação de categoria de POI e predição do próximo POI. Os resultados obtidos nos três estados avaliados indicam que sim. Ao incorporar um \textit{encoder} espacial contínuo, um \textit{encoder} temporal (Time2Vec) e um \textit{encoder} categórico hierárquico (HGI), o ST-MTLNet supera de forma consistente o \textit{baseline} baseado apenas em DGI.

Na tarefa de classificação de categoria de POI, as variantes propostas superam o MTLNet em todas as 21 combinações de categoria e estado, com ganhos médios de 20 a 24 pontos percentuais. Na tarefa de predição do próximo POI, o desempenho também é majoritariamente favorável ao ST-MTLNet, que vence em 16 das 21 combinações avaliadas, com destaque para categorias como \textit{Food}. Em conjunto, esses resultados mostram que a decomposição da representação em componentes espaciais, temporais e categóricos permite capturar padrões que o DGI, por si só, não modela de forma suficiente.

A comparação entre SIREN e Sphere2Vec-M também mostra que não há um único \textit{encoder} espacial universalmente superior. O SIREN apresenta melhor desempenho com maior frequência na Flórida e na Califórnia, enquanto o Sphere2Vec-M se destaca no Texas, sugerindo que a adequação do \textit{encoder} depende da distribuição geográfica dos POIs em cada território. Assim, o principal ganho do trabalho não está apenas na escolha de uma arquitetura espacial específica, mas na própria adoção de representações desacopladas como alternativa ao paradigma monolítico do DGI.

Uma limitação relevante permanece na categoria \textit{Travel} da tarefa de predição do próximo POI, na qual o MTLNet original ainda obtém os melhores resultados em parte dos cenários. Isso sugere que deslocamentos de longa distância e relações topológicas entre regiões ainda são melhor capturados pela estrutura de grafo utilizada pelo DGI. Além disso, como os três componentes propostos são utilizados em conjunto, este trabalho não isola a contribuição individual de cada \textit{encoder}, o que restringe uma análise mais detalhada sobre o peso relativo de cada fonte de informação.

Como trabalhos futuros, destacam-se a investigação de arquiteturas multitarefa mais flexíveis, a exploração de estratégias de fusão mais sofisticadas entre os \textit{encoders} e a combinação entre representações baseadas em grafo e codificações contínuas de coordenadas. Essas direções podem ampliar os ganhos observados e ajudar a reduzir as limitações ainda presentes em categorias como \textit{Travel}.